\section{Integration with the BX3 Framework}
\label{sec:integration}

The Recursive Spawning mechanism does not operate in isolation. Its correctness
guarantees are derived entirely from integration with the BX3 three-layer
architecture: the \textit{Purpose} layer serves as the accountability anchor,
the \textit{Bounds} layer enforces depth constraints, and the \textit{Fact} layer
maintains the forensic ledger. Each layer contributes a necessary and
sufficient condition for the upstream accountability guarantee. This section
specifies each integration point precisely.

% ---------------------------------------------------------------
\subsection{Purpose Layer as Accountability Anchor}
\label{sec:purpose-anchor}

Every node $v$ in a BX3 spawn tree carries a human anchor $h(v)$ that is
established at node provisioning and is structurally immutable for the node's
operational lifetime. The \textit{Purpose} layer is the architectural locus of
this anchor: it encodes the principal's intent, defines the operating range
$R(v)$, and receives all escalated exceptions that the \textit{Bounds} layer
cannot resolve.

The integration between Recursive Spawning and the \textit{Purpose} layer
operates at three distinct points:

\begin{enumerate}
\item[(a)] \textbf{Anchor inheritance at spawn.} When parent node $p$ spawns
child $c$, the spawn directive encodes $h(c) = h(p)$. This inheritance is not
a runtime policy decision — it is enforced by the \fact{} layer's pre-commit
hook before the child node is provisioned. The consequence is that every node
in the spawn tree, regardless of depth, points to the same human principal as
its ultimate accountability terminus.

\item[(b)] \textbf{Escalation terminus.} The Bailout Protocol traverses the
parent-pointer chain $\alpha(v), \alpha^{(2)}(v), \ldots$ until reaching
$h(v)$. Because $h(v)$ is always a human principal and never a machine actor,
the escalation path never terminates at an intermediate machine node. The
\textit{Purpose} layer at each node is the last machine layer before the human
anchor. This structural property makes the upstream accountability guarantee
non-vacuous: the guarantee holds precisely because the terminus is defined at
the \textit{Purpose} layer, which cannot be replicated or substituted by any
machine actor.

\item[(c)] \textbf{Scope adjudication.} When the \textit{Bounds} layer detects
a delegation scope inheritance violation, the exception is escalated to the
\textit{Purpose} layer of the parent node for adjudication. The
\textit{Purpose} layer alone has the authority to redefine $R(p)$ or to
authorize an exception to the inheritance constraint on behalf of the human
principal. No machine actor may grant this authorization.
\end{enumerate}

The combined effect of (a)–(c) is that the \textit{Purpose} layer serves as
a double anchor: it anchors the spawn tree's accountability chain at the top
(human principal $\rightarrow$ root node's \textit{Purpose} layer), and it
anchors every escalation path at its terminus. The spawn tree is not a tree of
autonomous agents — it is a tree of accountability structures, each anchored to
the same human principal.

% ---------------------------------------------------------------
\subsection{Bounds Engine: Depth Enforcement}
\label{sec:bounds-depth}

The \textit{Bounds} layer enforces the Hierarchy Finiteness Invariant
(Invariant~\ref{inv:hierarchy-finiteness}) as a deterministic pre-commit check
on every spawn event. The enforcement mechanism is a depth-evaluation function
$\textit{depth-ok}(c, p)$ that computes $d(c) = d(p) + 1$ and compares the result
against the deployment constant $D_{\max}$.

When a node in state \texttt{CAN\_SPAWN} receives an $e_{\text{spawn}}$ event,
it transitions to \texttt{SPAWN\_PENDING} and issues a depth-confirmation request
to the \textit{Bounds} layer. The \textit{Bounds} layer evaluates
$\textit{depth-ok}(c, p)$ and, if the check passes, emits $e_{\text{depth-ok}}$.
If the check fails, the \textit{Bounds} layer emits $e_{\text{depth-violation}}$,
transitions the spawning state machine to \texttt{SPAWN\_LOCKED}, triggers the
Bailout Protocol, and records the violation in the Forensic Ledger. The depth
check is a deterministic integer comparison — no probabilistic reasoning, no
model judgment. Either $d(c) \leq D_{\max}$ or it does not.

The critical property of this integration is that $D_{\max}$ is a deployment
parameter established at system initialization and is not negotiable at runtime.
No machine actor, including the \textit{Bounds} layer itself, may raise
$D_{\max}$ to accommodate a non-compliant spawn. The deployment designer sets
$D_{\max}$ once, based on the operational requirements of the deployment
context — specifically, the maximum acceptable escalation path length
$T_{\max}$ for the Bailout Protocol's Liveness guarantee. The \textit{Bounds}
layer's enforcement ensures that this contract is never violated.

The delegation scope inheritance check (Invariant~\ref{inv:scope-inheritance})
is also enforced by the \textit{Bounds} layer as a deterministic set comparison.
The check evaluates $R(c) \subset R(p)$ before every spawn event. If the
subset relation fails, the \textit{Bounds} layer emits $e_{\text{scope-violation}}$,
transitions the spawning state machine to \texttt{SPAWN\_LOCKED}, and triggers
the Bailout Protocol. The scope-inheritance check and the depth check are
evaluated in sequence (scope first, then depth), and both must pass before the
spawn proceeds to \texttt{CHILD\_PROVISIONED}. This sequencing ensures that the
first detected violation is the one recorded in the Forensic Ledger, which
matters for accountability attribution during human review.

% ---------------------------------------------------------------
\subsection{Fact Layer: Forensic Ledger}
\label{sec:fact-fledger}

The \fact{} layer provides two structural guarantees for recursive spawning:
immutable parent-pointer enforcement and complete, cryptographically chained
event recording.

\textbf{Immutable parent pointer.} The parent-pointer field $\alpha(v)$ is a
read-only property of the node's \fact{} layer worksheet. Any attempt to write
to $\alpha(v)$ at runtime is intercepted by the pre-commit hook and rejected.
The rejection event is recorded in the Forensic Ledger as a
\texttt{parent-pointer-write-denied} entry, and the Bailout Protocol is triggered
if the attempted write was initiated by a machine actor. This enforcement ensures
that the accountability chain is never redirected after node creation — the
chain that exists at provisioning is the chain that persists throughout the
node's operational lifetime.

\textbf{Cryptographically chained forensic record.} Every spawn event, every
state transition in the spawning state machine, every depth-confirmation
result, every scope-inheritance evaluation, and every human directive issued
by $h(v)$ in response to a spawn-related exception is recorded as a Forensic
Ledger entry. Each entry carries the SHA-256 hash of the preceding entry,
forming an append-only chain from the genesis entry to the most recent. The
cryptographic chaining ensures that no entry can be modified, deleted, or
inserted without breaking the chain — a condition detectable by any auditor
without access to the ledger's contents.

The Forensic Ledger entries relevant to recursive spawning include:
\texttt{spawn-proposed}$(p, c, R(c), d(c))$, \texttt{depth-confirmed},
\texttt{depth-rejected}, \texttt{scope-confirmed}, \texttt{scope-rejected},
\texttt{child-provisioned}$(p, c, \alpha(c)=p, h(c)=h(p))$,
\texttt{parent-pointer-write-denied}, and \texttt{spawn-locked}. This record is
the complete, immutable evidence base for reconstructing the accountability
chain of any node at any point in time. It is not a conventional audit log that
records actions after they are taken — it is an enforcement mechanism that
queries ledger state before subsequent actions are permitted, ensuring that the
ledger is the authoritative source of truth for the runtime state of the spawn
tree.

% ---------------------------------------------------------------
\subsection{Unified Integration Property}
\label{sec:unified-integration}

The three-layer integration yields a single structural property that is the
foundation of all downstream correctness guarantees:

\begin{invariant}[Spawn Tree Accountability Invariant]
\label{inv:spawn-accountability}
For every node $v$ in a BX3 spawn tree, the accountability chain
$\alpha^{(k)}(v) = h(v)$ with $k \leq D_{\max}$ is: (i) immutable — $\alpha(v)$
is never modified after node creation; (ii) complete — every spawn event and
state transition is recorded in the Forensic Ledger; (iii) terminating — the
chain terminates at a human principal $h(v)$, never at a machine actor; and
(iv) enforceable — any violation of (i), (ii), or (iii) triggers the Bailout
Protocol before the system state advances.
\end{invariant}

This invariant is not a policy constraint or an operational guideline. It is a
structural property enforced by the pre-commit hooks of all three layers
operating in concert. None of the three layers alone is sufficient: the
\textit{Purpose} layer provides the human anchor; the \textit{Bounds} layer
provides the depth and scope enforcement; the \textit{Fact} layer provides the
immutable parent pointer and the cryptographically chained forensic record.
Together they enforce Invariant~\ref{inv:spawn-accountability}, and it is this
invariant that enables the correctness properties established in the next
section.